\section{The Excavation: History Without Flattery}

\subsection{Episodes Over Adjectives}

Do not begin with adjectives such as \textit{analytical}, \textit{passionate}, \textit{resilient}, or \textit{innovative}. These are difficult to verify and apply to almost everyone. Begin with concrete operational episodes.

Take your last 5 to 10 years, or your entire career if shorter. Draw a timeline of jobs, academic projects, side ventures, volunteering, family businesses, failed initiatives, competitions, and independently pursued interests. \textbf{You are looking for repeated patterns of usefulness, not a flattering professional biography.}

For each period, answer:

\begin{itemize}[leftmargin=11pt,itemsep=1.5pt,topsep=1.5pt,parsep=0pt]
  \item What specific problem was I asked to solve?
  \item What did I personally do, separated from team contributions?
  \item What made the task difficult in practice?
  \item Who benefited, and how would they measure the improvement?
  \item What did I learn about customer incentives, pricing, and failure modes?
  \item What proof survives (code, reports, customer messages, logs, or references)?
  \item What did I discover that industry outsiders regularly misunderstand?
  \item If I had to repeat the work without the employer's brand, staff, or capital, what could I still produce?
\end{itemize}

That last question is a rigorous filter. An executive may have managed a unit generating hundreds of crores, but the results depended on infrastructure and brand equity they cannot export to a new firm.

\subsection{4 Less Obvious Signals}

\begin{itemize}[leftmargin=11pt,itemsep=2pt,topsep=1.5pt,parsep=0pt]
  \item \textbf{The repeated request:} What do colleagues routinely ask you to diagnose? Explain difficult concepts? Debug production outages? Evaluate job candidates? Price complex proposals? These requests reveal perceived competence.
  \item \textbf{The operational rescue:} When projects were failing or ambiguous, what were you trusted to repair? Crises reveal tacit judgment, calm triage, and cross-functional negotiation.
  \item \textbf{The chronic frustration:} Which recurring industry workflows angered you because you saw how needlessly broken they were? Frustration points to gaps, but check whether buyers care enough to pay.
  \item \textbf{The authentic obsession:} What problems did you study without external deadlines? Sustained curiosity matters because early ventures require endless unglamorous execution.
\end{itemize}

\subsection{The 10-Episode Exercise}

Select 10 concrete episodes. Write each in 5 sentences:

\begin{enumerate}[label=\textbf{\arabic*.},leftmargin=11pt,itemsep=1.5pt,topsep=1.5pt,parsep=0pt]
  \item \textit{Situation:} Where was I, who needed something, and why?
  \item \textit{Action:} What did I personally do?
  \item \textit{Outcome:} What changed, with numbers if available?
  \item \textit{Mechanism:} Which capability or insight made the result possible?
  \item \textit{Portability:} Could I produce the same result for an independent client today?
\end{enumerate}

Consider an example:

\begin{tcolorbox}[
  enhanced,
  colback=BayesPale!25!white,
  colframe=BayesBorder!60!white,
  boxrule=0.4pt,
  arc=2pt,
  left=6pt,right=6pt,top=4pt,bottom=4pt,
  before skip=3pt,after skip=4pt
]
\footnotesize\color{BayesInk}
An operations manager repeatedly faced unexplained stockouts. I combined sales history with supplier lead-time data, built a simple replenishment rule, and coached the planning team. Over the pilot, emergency orders fell, although the evidence is too limited to attribute the entire improvement to my intervention. My useful capability was turning operational ambiguity into a quantitative workflow people could follow. To commercialize it independently, I would still need buyer access, permission to use appropriate data, and proof that the approach works in another company.
\end{tcolorbox}

\textbf{Notice the difference between describing an impressive project and identifying a portable capability: the latter is where authentic business hypotheses originate.}

\subsection{The Personal Advantage Ledger}

Create a ledger and record verifiable facts rather than personality narratives:

\begin{tcolorbox}[
  enhanced,
  colback=white,
  colframe=BayesBorder!80!white,
  boxrule=0.4pt,
  arc=2pt,
  left=3pt,right=3pt,top=3pt,bottom=3pt,
  before skip=3pt,after skip=5pt
]
\scriptsize
\setlength{\tabcolsep}{2.5pt}
\renewcommand{\arraystretch}{1.18}
\begin{tabularx}{\linewidth}{@{} >{\bfseries\color{BayesInk}}p{2.3cm} >{\color{BayesInk}}p{2.2cm} >{\color{BayesSecondary}}X @{}}
\toprule
\textbf{Hypothesis \& Artifact} & \textbf{Target Buyer} & \textbf{Key Uncertainty} \\
\midrule
Teach technical ideas (recordings, student tests) & Developers, engineering training teams & Can I deliver reliably without personal presence? \\
\midrule
Diagnose inventory waste (documented workflows) & Mid-size parts manufacturers & Will owners pay external fees for pure diagnosis? \\
\midrule
Deploy reliable software (production code, post-mortems) & Engineering teams with tight SLAs & Can I sell and support systems as a sole founder? \\
\bottomrule
\end{tabularx}
\end{tcolorbox}

Keep weak evidence visible. You will return to this ledger repeatedly as your venture matures.
